\section{Introduction}

The rapid evolution of AI coding agents has shifted their operational paradigm from passive code-completion suggestions to autonomous systems capable of executing real actions on live repositories. While this enhances developer productivity, it introduces critical security risks. Unlike traditional static analysis tools, autonomous agents interact with the environment, execute shell commands, and modify codebases iteratively, making their execution-time security posture a paramount concern.

This study aims to systematically evaluate the execution security of state-of-the-art coding agents. Our scope focuses on four primary attack vectors: Prompt Injection (PI), Context Poisoning (CP), Dependency Confusion (DC), and Secret Leakage (SL). These vectors were selected from a candidate list of seven because they present the most immediate, measurable threats that can be consistently reproduced across different agent architectures. We deliberately excluded Retrieval Poisoning, Tool Abuse, and Memory Poisoning from this study. Evaluating these excluded vectors would have compromised methodological rigor due to the significant architectural disparities among the agents (e.g., chat-based interfaces vs. CLI-driven tools), which prevents cross-agent comparability for those specific vulnerabilities. By focusing on the four selected vectors, we establish a robust, comparable benchmark that highlights how different delivery mechanisms and agent architectures impact susceptibility to execution-time exploitation.
